% registry: WANG2018
\methodhead{WANG2018}{Wang, Shang, Sun and Chen (2018): weighted critical plane with mean shear stress}\label{sheet:WANG2018}

\block{Source} X.-W.~Wang, D.-G.~Shang, Y.-J.~Sun, H.~Chen, \emph{Theor.\ Appl.\ Fract.\ Mech.}
(2018), accepted manuscript \cite{wang2018}, Eqs.~(1)--(20); Tables~1--5; Section~5.1.

\block{Idea and scope} The critical plane is the weighted time average of the planes of maximum
shear stress. The weights grow with the instantaneous maximum shear stress. On that plane, the
shear and normal stress amplitudes are combined in von Mises form. Each is corrected by its own
mean stress: the shear mean through its absolute value, the normal mean with its sign. The life is
read on the fully reversed torsion S--N curve. Proposed for high-cycle lives of 7075-T651
thin-walled tubes under tension--torsion: constant amplitude, in and out of phase, with mean
stresses, and variable amplitude.

\block{Inputs}
\begin{center}\small
\begin{tabular}{@{}llp{100mm}@{}}
\toprule
Symbol & Unit & Meaning\\\midrule
$\sigma(t)$, $\tau(t)$ & MPa & axial normal stress and shear stress of the tube ($\sigma_{yy}$, $\tau_{xy}$ in the repository frame)\\
$\varphi$ & rad & angle between the normal of a plane perpendicular to the surface and the axis\\
$\sigma_\varphi$, $\tau_\varphi$ & MPa & normal and shear stress on plane $\varphi$, Eq.~\eqref{eq:w18-plane}\\
$\tau_{\max}(t_k)$ & MPa & maximum shear stress over the surface planes at instant $t_k$\\
$\varphi_{45}(t_k)$ & rad & orientation of the maximum-shear plane at $t_k$, in $[-45^\circ,45^\circ]$\\
$W(t_k)$ & -- & weight of instant $t_k$, Eq.~\eqref{eq:w18-weight}\\
$\hat\varphi$, $\varphi_c$ & rad & weighted maximum-shear plane and critical plane\\
$\tau_a,\tau_m$; $\sigma_a,\sigma_m$ & MPa & amplitude and mean of $\tau_\varphi$ and $\sigma_\varphi$ on $\varphi_c$\\
$\tau_{-1}=\tm$, $\sigma_{-1}=\fm$ & MPa & fully reversed torsion and axial limits\\
$\sigma_0$, $\tau_0$ & MPa & pulsating ($R=0$) axial and torsion limits, maximum stress\\
$k_\tau$, $k_\sigma$ & -- & mean shear and mean normal stress sensitivities, Eq.~\eqref{eq:w18-k}\\
$\tau_{eq}$ & MPa & equivalent shear stress, Eq.~\eqref{eq:w18}\\
\bottomrule
\end{tabular}
\end{center}

\block{Formulation} On the planes perpendicular to the surface (Eq.~2),
\begin{equation}
\sigma_\varphi = \tfrac{\sigma}{2}+\tfrac{\sigma}{2}\cos2\varphi+\tau\sin2\varphi ,\qquad
\tau_\varphi = \tfrac{\sigma}{2}\sin2\varphi-\tau\cos2\varphi .
\label{eq:w18-plane}
\end{equation}
At each instant $t_k$ of one cycle (Eqs.~3--7),
\begin{gather}
\tau_{\max}(t_k)=\sqrt{\bigl(\sigma(t_k)/2\bigr)^2+\tau(t_k)^2},\qquad
\varphi_{45}(t_k)=\tfrac12\arctan\!\left[-\frac{\sigma(t_k)}{2\tau(t_k)}\right],\\
W(t_k)=\frac{\tau_{\max}(t_k)-\min_k\tau_{\max}}{\max_k\tau_{\max}-\min_k\tau_{\max}},\qquad
\lvert\hat\varphi\rvert=\frac{\sum_k\lvert\varphi_{45}(t_k)\rvert\,W(t_k)}{\sum_k W(t_k)} ,
\label{eq:w18-weight}
\end{gather}
and $\hat\varphi$ takes the sign of $\sum_k\varphi_{45}(t_k)W(t_k)$. The shear stress has the same
magnitude on $\hat\varphi$ and $\hat\varphi+90^\circ$, so the candidate with the larger damage is
$\varphi_c$. When $\tau_{\max}$ is constant (Eqs.~8--9), the plane of the largest normal stress is
critical. On $\varphi_c$ (Eqs.~11, 15, 18),
\begin{gather}
\tau_{eq}=\sqrt{\bigl[\tau_a+k_\tau\lvert\tau_m\rvert\bigr]^2+\tfrac13\bigl[\sigma_a+k_\sigma\sigma_m\bigr]^2}\le\tau_{-1},
\label{eq:w18}\\
k_\tau=\frac{2\tau_{-1}}{\tau_0}-1,\qquad
k_\sigma=\sqrt{\frac{48\tau_{-1}^2}{\sigma_0^2}-\frac{12\tau_{-1}^2}{\tau_0^2}}-1 ,
\label{eq:w18-k}
\end{gather}
obtained by imposing reversed torsion ($\tau_a=\tau_{-1}$), pulsating tension
($\tau_a=\tau_m=\sigma_a=\sigma_m=\sigma_0/4$ on the $45^\circ$ plane) and pulsating torsion
($\tau_a=\tau_m=\tau_0/2$), Eqs.~(12)--(14). Eq.~(17) requires $k_\tau>0$ and $k_\sigma>0$.
Without the pulsating limits (Eq.~19),
\begin{equation}
\sigma_0=\sqrt2\,\sigma_{-1},\qquad \tau_0=\frac{4\tau_{-1}}{2\sigma_{-1}/\sigma_0+1},
\end{equation}
which gives $k_\tau=(\sqrt2-1)/2$ (Eq.~20). Reversed tension gives $\tau_{eq}=\sigma_a/\sqrt3$: the
criterion implies $\sigma_{-1}=\sqrt3\,\tau_{-1}$.

\block{Constants and calibration} $k_\tau$, $k_\sigma$ from $\tau_{-1}$, $\sigma_0$, $\tau_0$, or
from $\sigma_{-1}$, $\tau_{-1}$ through Eq.~(19). The paper takes $\sigma_{-1}$ and $\tau_{-1}$
from its S--N curves at $N_0=10^6$, $2\cdot10^6$ and $10^7$, and compares the three choices.

\block{Life} $\tau_{eq}$ on the fully reversed torsion S--N curve (Section~4.2). Variable
amplitude: Wang--Brown reversals on the two candidate planes, Miner's rule, and the larger damage
(Eqs.~21--24).

\block{Repository choices}
\begin{itemize}
\item Scope: tension--torsion only (maintainer decision). Any other stress component is out of
  scope.
\item One cycle sampled at 3600 instants for Eqs.~\eqref{eq:w18-weight}. On a plane, the
  amplitude and mean are exact single-harmonic values.
\item The two candidates are compared through $\tau_{eq}$: there is one S--N curve, so the larger
  $\tau_{eq}$ is the larger damage. A zero signed sum gives $\hat\varphi\ge0$.
\item Constant $\tau_{\max}$: the plane of the largest $\max_t\sigma_\varphi(t)$ is found over
  3600 planes.
\item $\sigma_{-1},\tau_{-1}$ are the workbook \texttt{Sigm1}, \texttt{Taum1}. $\sigma_0,\tau_0$
  are \texttt{Sig0}, \texttt{Tau0}, otherwise Eq.~(19). A violation of Eq.~(17) is out of material
  scope.
\item Life: the master curve $\tau_{eq}=AN^b$ is fitted on the $R=-1$ pure torsion failures,
  which is the paper's torsion S--N curve. The variable-amplitude procedure is not implemented.
\end{itemize}

\block{Worked example} 7075-T651, Table~2 curves at $N_0=10^6$: $\sigma_{-1}=222.0$,
$\tau_{-1}=127.1\MPa$, so $k_\tau=0.2071$ and $k_\sigma=0.8698$. Table~5 load
$\sigma=163.02+144.57\sin\omega t$, $\tau=94.12+83.47\sin(\omega t-\delta)$ (MPa):
\begin{center}\small
\begin{tabular}{@{}lccc@{}}
\toprule
$\delta$ & $\varphi_c$ & $\tau_{eq}$ (MPa) & $N$ from $\tau=688.85N^{-0.12233}$\\\midrule
$0^\circ$ & $69.6^\circ$ & 159.3 & 157\,636\\
$30^\circ$ & $-20.4^\circ$ & 160.0 & 152\,280\\
$90^\circ$ & $-20.1^\circ$ & 160.4 & 149\,120\\
\bottomrule
\end{tabular}
\end{center}
Experimental lives: 53\,058--65\,046, 64\,231--82\,445 and 70\,333--94\,718 cycles.

\block{Verification} Section~5.1: the estimated crack trace, $90^\circ-\lvert\varphi_c\rvert$, is
$70^\circ$ for every $30^\circ$ and $90^\circ$ out-of-phase test of Table~5, as in the paper.
Eqs.~(12)--(14) give $\tau_{eq}=\tau_{-1}$ exactly; Eq.~(20) equals Eq.~(15) with Eq.~(19). The
paper prints the predicted lives only in figures.
